\documentclass[11pt]{article}

\usepackage[T1]{fontenc}
\usepackage[utf8]{inputenc}
\usepackage[margin=1.04in]{geometry}
\usepackage{amsmath,amssymb,amsthm,mathtools,bm}
\usepackage{booktabs,array,tabularx,enumitem,microtype}
\usepackage[colorlinks=true,linkcolor=blue,citecolor=blue,urlcolor=blue]{hyperref}
\usepackage[nameinlink,capitalise]{cleveref}
\setlength{\emergencystretch}{3em}
\setlength{\parskip}{0.16em}
\setlength{\parindent}{1.35em}
\raggedbottom
\numberwithin{equation}{section}

\newtheorem{theorem}{Theorem}[section]
\newtheorem{proposition}[theorem]{Proposition}
\newtheorem{lemma}[theorem]{Lemma}
\newtheorem{corollary}[theorem]{Corollary}
\newtheorem{conjecture}[theorem]{Conjecture}
\newtheorem{assumption}[theorem]{Assumption}
\theoremstyle{definition}
\newtheorem{definition}[theorem]{Definition}
\newtheorem{principle}[theorem]{Principle}
\newtheorem{interpretation}[theorem]{Interpretive Thesis}
\theoremstyle{remark}
\newtheorem{remark}[theorem]{Remark}
\newtheorem{openproblem}[theorem]{Open Problem}

\newcommand{\C}{\mathbb C}
\newcommand{\R}{\mathbb R}
\newcommand{\Z}{\mathbb Z}
\newcommand{\Hilb}{\mathcal H}
\newcommand{\B}{\mathcal B}
\newcommand{\Aalg}{\mathcal A}
\newcommand{\Ralg}{\mathcal R}
\newcommand{\Oalg}{\mathcal O}
\newcommand{\Qpred}{\mathsf Q_{\rm pred}}
\newcommand{\Instr}{\mathcal I}
\newcommand{\Rec}{\mathsf R}
\newcommand{\Tr}{\operatorname{Tr}}
\newcommand{\supp}{\operatorname{supp}}
\newcommand{\Id}{\mathbf 1}
\newcommand{\Kfour}{K_4}
\newcommand{\Hgen}{\mathcal H_{\rm gen}}
\newcommand{\Mhor}{\mathcal M_{\rm hor}}
\newcommand{\Mrad}{\mathcal M_{\rm rad}}
\newcommand{\SBH}{S_{\rm BH}}
\newcommand{\Srad}{S_{\rm rad}}
\newcommand{\CP}{\operatorname{CP}}

\title{\textbf{Renewal Records and Apparent Wavefunction Collapse}\\[2mm]
\large Predictive Conditioning, Born Weights, and the Horizon Record Test}
\author{A.~P\'elissier}
\date{}

\begin{document}
\maketitle

\begin{abstract}
We formulate a renewal interpretation of wavefunction collapse.  The proposal is not an objective-collapse dynamics: the Schr\"odinger equation is not supplemented by stochastic localization, and decohered branches are not postulated to be dynamically annihilated.  Collapse is instead identified with predictive conditioning on a stable regenerative record.  A measurement outcome is a sector of the observer's record algebra, and the L\"uders rule is the conditional state induced on that sector.  This places the renewal account near decoherence, consistent histories, relational quantum mechanics, and Everett-compatible unitary evolution, while giving a sharper criterion for effective branching: a branch is not an arbitrary component of a state-vector expansion, but a stable sector of a regenerative predictive quotient.  The Born rule is correspondingly recast as a disintegration theorem for stable records.  Pointer-sector Born weights first follow from additivity on the Boolean algebra generated by stable records.  The later renewal Born argument then upgrades this statement: additivity follows from predictive refinement stability, multiplicativity from independent renewal composition, and phase/deck invariance from channel equivalence, forcing quadratic branch weights.  A projective two-design closure of the reset ensemble extends the rule from pointer sectors to all implementable measurement bases.  The black-hole information problem enters only as a stringent application of the same logic.  In a horizon setting, Hawking radiation is an exterior record stream, and Page recovery asks whether that stream becomes complete enough to reconstruct the branch data that crossed the horizon.  Thus measurement and black-hole information are not treated as two unrelated puzzles: ordinary measurement concerns branch selection by records, while horizon evaporation concerns branch reconstruction from records.  The common object is the regenerative record algebra.
\end{abstract}

{\setlength{\parskip}{0pt}\tableofcontents}

\section{Introduction}
\label{sec:intro}

The measurement problem is often presented as a conflict between two descriptions of quantum evolution.  On the one hand, the microscopic state evolves linearly and preserves superpositions.  On the other hand, laboratory observations are definite: a detector clicks in one channel, a pointer occupies one macroscopic position, and an observer remembers one outcome.  The tension is not merely technical.  It concerns the relation between the mathematical state, the physical record, and the ontology assigned to alternatives that no longer interfere.

Decoherence explains why macroscopic alternatives rapidly cease to interfere when information about them is dispersed into an environment \cite{Zeh1970,Zurek1981,JoosZeh1985,Schlosshauer2007}.  It does not, by itself, decide what should be said about the alternatives after interference becomes inaccessible.  Everettian interpretations retain the global state and read the decohered components as branches of the universal wavefunction \cite{Everett1957,DeWittGraham1973,Wallace2012}.  Objective-collapse theories add a real stochastic modification of the dynamics \cite{GRW1986,Pearle1989,BassiGhirardi2003}.  Operational or Copenhagen views treat collapse as a rule for updating predictions.  Consistent-histories approaches attach probabilities to decoherent families of histories \cite{Griffiths1984,Omnes1992,GellMannHartle1993}.  Relational and QBist interpretations make state assignment dependent on the system, observer, or agent relative to which information is being encoded \cite{Rovelli1996,FuchsMerminSchack2014}.  Each view locates definiteness in a different place: in dynamics, in branches, in histories, in relations, or in agents.

The renewal proposal developed here locates definiteness in a record algebra.  A quantum state is not taken to be an observationally complete object until the algebra of stable records relative to which predictions are made has also been specified.  In this sense collapse is neither a new microscopic force nor a merely psychological event.  It is the passage from an unconditioned predictive state to the conditional state determined by a stable regenerative record.  Symbolically,
\[
  \boxed{\text{apparent collapse} = \text{conditioning on a stable record sector}.}
\]
The word ``apparent'' is used in a precise sense.  It denies a fundamental nonunitary jump in the underlying dynamics; it does not deny the reality of the outcome for observers whose future predictions are now conditioned on a record.

The mathematical language is inherited from renewal geometry.  Microscopic histories are not distinguished by all their microscopic details, but by the accumulated completely positive channel they induce.  Histories that induce the same future predictive channel are identified in the predictive quotient.  A measurement record is then a stable central sector of the accessible algebra that survives this quotient and organizes future expectations.  This gives the usual measurement update
\[
  \rho \longmapsto \frac{P_i\rho P_i}{\Tr(P_i\rho)}
\]
a different interpretation: it is not the trace of a physical destruction of all other branches, but the conditional state in the sector selected by the retained record.

The same distinction also clarifies the relation to black holes.  A separate renewal black-hole theory develops horizon response, non-Markovian evaporation, Page timing, and the loss of controlled interior reconstruction.  The present paper does not reproduce that analysis.  It uses the horizon only as a conceptual stress test for measurement theory.  In an ordinary laboratory, records select the branch experienced by the observer.  In an evaporating black hole, exterior radiation must become a sufficiently complete record of the branch data that crossed the horizon.  Measurement asks how a record selects an outcome; black-hole information asks whether a record stream can reconstruct what is otherwise hidden.  The record algebra is the common object.

The paper proceeds as follows.  \Cref{sec:renewal-records} develops the renewal notion of a regenerative record and explains why it replaces a primitive classical cut.  \Cref{sec:collapse} formulates apparent collapse as conditionalization on a central record sector and proves the corresponding no-signalling statement.  \Cref{sec:born} treats pointer states, branch weights, Born disintegration, and objectivity.  \Cref{sec:interpretations} compares the resulting position with the principal interpretations of quantum mechanics.  \Cref{sec:phenomenology} isolates the empirical target as record-stream statistics rather than stochastic localization noise.  \Cref{sec:horizon} explains how black-hole information enters as the horizon record-completeness problem.  \Cref{sec:outlook} summarizes the metaphysical and technical status of the proposal.

\section{Regenerative records and the predictive quotient}
\label{sec:renewal-records}

The renewal framework begins not from a preferred classical apparatus, but from a family of update channels.  Let a microscopic history \(\omega\) determine a completely positive map \(\Phi_\omega\) on an appropriate algebra of predictive states.  The predictive quotient identifies histories that have the same accumulated effect:
\begin{equation}
\label{eq:predictive-quotient}
  \omega\sim\omega'
  \quad\Longleftrightarrow\quad
  \Phi_\omega=\Phi_{\omega'} .
\end{equation}
This quotient expresses a simple operational principle: microscopic distinctions that never change any future prediction in the accessible algebra should not be treated as distinct effective alternatives.

A measurement record is a special kind of surviving distinction.  It is not merely a correlation between a system and an environment.  It is a correlation that becomes stable under future updates and can therefore serve as a conditioning datum for later predictions.  A scattered photon, a pointer position, a detector discharge, and a human memory all become records only insofar as they retain such predictive force.

\begin{definition}[Record algebra]
Let \(\Hilb\) be a Hilbert space containing the system, apparatus, and environment, and let \(\Aalg\subseteq\B(\Hilb)\) be the algebra accessible to an observer or macroscopic subsystem.  A \emph{record algebra} is a subalgebra \(\Ralg\subseteq\Aalg\) generated by effects \(\{E_i\}\) whose outcome labels are stable under subsequent renewal and whose distinctions survive the predictive quotient \eqref{eq:predictive-quotient}.
\end{definition}

The emphasis on stability is essential.  A transient microscopic correlation may temporarily encode information, but it is not a measurement outcome unless the future dynamics treats it as a retained causal state.  In renewal language, the apparatus must regenerate through the record.

\begin{definition}[Regenerative record]
A record sector \(i\) is \emph{regenerative} if there is a completely positive projection \(\Pi_i\) such that, after the record is written, future predictive evolution factors through the sector,
\begin{equation}
\label{eq:regenerative-factorization}
  \Phi_{\rm future}\circ\Phi_{\rm record}
  \simeq
  \Phi_{\rm future}^{(i)}\circ\Pi_i,
\end{equation}
up to errors that are invisible in the predictive quotient or decay under repeated renewal.
\end{definition}

This definition turns the usual vague language of ``macroscopic definiteness'' into a structural condition.  A definite outcome is a stable sector of the future-predictive algebra.  The observer's experience is definite because the observer is physically embedded in one such sector, not because a bare state vector has acquired an additional metaphysical property.

\subsection{Classical sectors from entanglement-breaking records}

The cleanest record limit is the entanglement-breaking limit.  A measurement instrument has completely positive components
\begin{equation}
\label{eq:instrument-components}
  \Instr_i(\rho)=K_i\rho K_i^*,
  \qquad
  p_i=\Tr\Instr_i(\rho),
  \qquad
  \rho_i=\frac{\Instr_i(\rho)}{p_i},
\end{equation}
and a nonselective channel \(\Instr=\sum_i\Instr_i\).  The record is effectively classical when the observer's algebra contains the outcome labels but no longer contains phase-sensitive operators connecting distinct labels.

\begin{assumption}[Predictively sufficient pointer record]
\label{ass:pointer-record}
A family of record effects \(\{E_i\}\) is a predictively sufficient pointer record when:
\begin{enumerate}[label=\textup{(\roman*)},leftmargin=2.2em]
\item the outcome labels are distinguishable in \(\Ralg\);
\item off-sector interference operators are not retained in the predictive quotient;
\item the same sector labels are preserved under repeated renewal; and
\item branch weights are additive under disjoint coarse graining of sectors.
\end{enumerate}
\end{assumption}

This assumption replaces the classical/quantum cut.  The cut is not a primitive boundary between two kinds of matter.  It is the emergent boundary of a record algebra whose sectors are stable enough to support conditional prediction.

\section{Apparent collapse as conditionalization}
\label{sec:collapse}

The textbook collapse rule can be written for a projective record as
\begin{equation}
\label{eq:luders}
  \rho\longmapsto \rho_i=\frac{P_i\rho P_i}{\Tr(P_i\rho)}.
\end{equation}
Objective-collapse theories interpret \eqref{eq:luders} as the effective trace of a real physical process in which the state is dynamically localized.  The renewal interpretation keeps the formula but changes its meaning.  The transition is not a new term in the microscopic equation of motion; it is the conditional state induced after the record sector \(i\) has entered the predictive algebra.

\begin{principle}[Renewal collapse principle]
\label{principle:collapse}
A collapse event is the transition from an unconditioned predictive state to the conditional state determined by a stable regenerative record:
\begin{equation}
\label{eq:renewal-collapse}
  \rho\big|_{\Aalg}
  \quad\longrightarrow\quad
  \rho\big|_{\Aalg,i}
  =\frac{\Pi_i(\rho)}{\Tr\Pi_i(\rho)}.
\end{equation}
It is physically irreversible for the observer when the sector \(i\) is stable under future renewal and redundantly encoded in the accessible algebra.
\end{principle}

The principle does not require the global state to cease evolving unitarily.  If the larger algebra contains all environmental phases, the global state may remain pure.  If the observer's algebra contains only the central record labels, the observer's state is sector-conditioned.  The apparent discontinuity is therefore algebra-relative: it is discontinuous as an update of predictive expectations, not necessarily as a microscopic dynamical evolution.

\subsection{Definiteness and sector stability}

Consider the standard premeasurement form
\begin{equation}
\label{eq:measurement-entanglement}
  \left(\sum_i c_i |s_i\rangle\right)|A_0\rangle|E_0\rangle
  \longrightarrow
  \sum_i c_i |s_i\rangle |A_i\rangle |E_i\rangle,
\end{equation}
with \(\langle E_i|E_j\rangle\simeq0\) for \(i\ne j\).  Decoherence explains the suppression of interference terms in local observables.  Renewal adds a criterion for when the labels in \eqref{eq:measurement-entanglement} count as outcomes: the apparatus-environment states \(|A_i\rangle|E_i\rangle\) must define regenerative sectors of the future-predictive algebra.

\begin{proposition}[Definiteness from central record sectors]
\label{prop:definiteness}
Let \(\{Z_i\}\) be orthogonal central projections in the observer's predictive algebra, with \(\sum_iZ_i=1\).  Suppose all future retained observables commute with the center generated by the \(Z_i\).  Then an observer whose record lies in sector \(i\) assigns the conditional state
\begin{equation}
\label{eq:central-conditioning}
  \rho_i=\frac{Z_i\rho Z_i}{\Tr(Z_i\rho)}.
\end{equation}
No observable retained in the predictive quotient detects coherence between distinct sectors.
\end{proposition}

\begin{proof}
Every retained observable decomposes as \(A=\sum_i Z_iAZ_i\).  Hence
\[
  \Tr(\rho A)=\sum_i \Tr(Z_i\rho Z_i A).
\]
Once the record label \(i\) is included in the observer's algebra, future expectations are computed with \eqref{eq:central-conditioning}.  Cross terms \(Z_i\rho Z_j\), \(i\ne j\), have no matrix element against any retained future observable.  They may be present in a larger algebra, but they are not predictive distinctions for the sector observer.
\end{proof}

This result captures the intended sense in which collapse is real without being fundamental.  It is real as a transition in the algebra of future predictions.  It is not fundamental as a universal destruction of all non-selected components.

\subsection{Locality and no signalling}

A common worry about collapse is that it appears instantaneous.  If collapse is interpreted as a physical process propagating through space, it threatens relativistic causality.  In the renewal account the update is instead conditionalization on an accessible record.  The standard no-signalling theorem is therefore not an additional constraint but a consistency check of the interpretation.

\begin{proposition}[No-signalling]
\label{prop:no-signalling}
Let \(\rho_{AB}\) be a bipartite state and let Alice's local record instrument satisfy \(\sum_i\Instr_i^A\) trace preserving.  Bob's nonselective reduced state is unchanged:
\begin{equation}
\label{eq:no-signalling}
  \rho_B'
  =\Tr_A\left[\sum_i(\Instr_i^A\otimes{\rm id}_B)(\rho_{AB})\right]
  =\Tr_A\rho_{AB}.
\end{equation}
Bob can condition on Alice's sector only after the label \(i\) enters Bob's accessible record algebra.
\end{proposition}

\begin{proof}
The identity is the usual trace-preserving property of a local quantum operation under the partial trace.  The sector-conditioned state depends on the classical label \(i\), and that label is not present in Bob's accessible algebra until it is physically communicated or otherwise recorded there.
\end{proof}

The proposition makes explicit why renewal collapse is compatible with relativistic causality.  Collapse is an update of the state relative to a record algebra, not a signal sent by the wavefunction.

\section{Born weights, pointer states, and objectivity}
\label{sec:born}

The preceding sections explain how definite sectors arise.  They do not yet explain why the weights of those sectors are the Born weights.  This is the deepest part of any interpretation that does not insert collapse probabilities as a primitive law.  The renewal strategy is to attach probability only to stable record sectors, then derive the algebraic properties of their weights from the predictive quotient itself.  The pointer-sector theorem is the first step.  The stronger Born-disintegration theorem below shows that, once record weights are stable under refinement, independent composition, and channel-equivalent phase/deck transformations, the quadratic rule is forced.

\subsection{Pointer-sector Born weights}

\begin{definition}[Predictive branch weight]
For a stable record sector \(i\), the branch weight is
\begin{equation}
\label{eq:branch-weight}
  w(i)=\Tr(E_i\rho),
\end{equation}
where \(E_i\) is the effect corresponding to the sector.  The weight is predictive if it is invariant under microscopic refinements of the history that leave the accumulated channel unchanged.
\end{definition}

\begin{theorem}[Pointer-sector Born rule]
\label{thm:pointer-born}
Assume \cref{ass:pointer-record}.  Then the unique normalized additive probability measure on the Boolean algebra generated by the pointer records is
\begin{equation}
\label{eq:pointer-born}
  p_i=\Tr(E_i\rho).
\end{equation}
For projective pointer records \(E_i=P_i\), the corresponding conditional state is the L\"uders state \eqref{eq:luders}.
\end{theorem}

\begin{proof}
The stable record sectors generate a commutative algebra.  Positivity, normalization, and additivity under disjoint coarse graining define a probability measure on this algebra.  The restriction of the quantum state \(\rho\) to the commutative algebra generated by the effects \(E_i\) gives precisely the measure \(p_i=\Tr(E_i\rho)\).  If the effects are orthogonal projectors, conditioning on the sector gives the L\"uders state.
\end{proof}

The theorem is deliberately the pointer-algebra part of a larger statement.  It assumes additivity on the Boolean algebra of stable records.  In the renewal formulation this additivity is not an independent probabilistic postulate once record sectors are required to be invariant under predictive refinement.  If a sector is split into two mutually orthogonal stable subrecords and the coarse-grained channel is unchanged, the weight of the coarse record must be the sum of the weights of the refined records.  Otherwise two predictively equivalent instruments would give different future statistics.

\subsection{Born disintegration from renewal multiplicativity}
\label{subsec:born-disintegration}

The full Born rule requires more than a commutative pointer algebra.  It requires that the weight assigned to an implemented record sector be independent of the microscopic representation of the same predictive channel, compatible with independent composition, and insensitive to unobservable phase or deck transformations.  These are structural properties of the predictive quotient rather than additional collapse dynamics.  The formulation below is the record-algebra version of the renewal Born-disintegration theorem \cite{PelissierBorn}.

\begin{definition}[Regular predictive amplitude weight]
\label{def:regular-weight}
Let \(\Hilb_{\rm pred}\) be the predictive Hilbert module obtained after quotienting histories by their accumulated completely positive channel.  A branch-weight functional \(W:\Hilb_{\rm pred}\to\R_{\ge0}\) is a \emph{regular predictive amplitude weight} when:
\begin{enumerate}[label=\textup{(\roman*)},leftmargin=2.2em]
\item \(W(0)=0\) and \(W(\psi)>0\) for nonzero retained predictive amplitudes;
\item if \(\psi\perp\phi\) are amplitudes in disjoint stable record sectors, then
\begin{equation}
\label{eq:weight-additivity}
  W(\psi+\phi)=W(\psi)+W(\phi);
\end{equation}
\item if two independently renewed systems are prepared with predictive amplitudes \(\psi\) and \(\phi\), then
\begin{equation}
\label{eq:weight-multiplicativity}
  W(\psi\otimes\phi)=W(\psi)W(\phi);
\end{equation}
\item if \(u\) is a phase or deck transformation that leaves the accumulated predictive channel unchanged, then \(W(u\psi)=W(\psi)\);
\item \(W\) is weakly regular, for example measurable or continuous on one nontrivial ray.
\end{enumerate}
\end{definition}

The assumptions have direct record-theoretic meanings.  Additivity expresses invariance under stable refinement and coarse graining of record sectors.  Multiplicativity expresses the tensor product of independent renewal experiments.  Phase/deck invariance expresses the defining equivalence relation of the predictive quotient: transformations that do not alter any future retained channel cannot alter the branch weight.

\begin{theorem}[Renewal Born disintegration]
\label{thm:renewal-born-disintegration}
Let \(W\) be a regular predictive amplitude weight on \(\Hilb_{\rm pred}\).  Then, after normalization on unit predictive amplitudes,
\begin{equation}
\label{eq:quadratic-weight}
  W(\psi)=\|\psi\|_{\rm pred}^{2}.
\end{equation}
Consequently, for an implemented projective record instrument \(\{\Pi_i\}\) and a pure predictive state \(\psi\),
\begin{equation}
\label{eq:born-disintegration}
  p_i
  =
  \frac{\|\Pi_i\psi\|_{\rm pred}^{2}}
    {\sum_j\|\Pi_j\psi\|_{\rm pred}^{2}}.
\end{equation}
For a density operator \(\rho\), the corresponding expression is
\begin{equation}
\label{eq:density-born}
  p_i=\Tr(\Pi_i\rho),
  \qquad
  \rho_i=\frac{\Pi_i\rho\Pi_i}{\Tr(\Pi_i\rho)} .
\end{equation}
\end{theorem}

\begin{proof}
By phase and deck invariance, the weight of a one-dimensional retained amplitude depends only on its predictive norm.  Write \(W(\psi)=F(\|\psi\|_{\rm pred})\) on rays, and set \(G(x)=F(\sqrt{x})\).  If \(\psi\perp\phi\), then
\[
  \|\psi+\phi\|_{\rm pred}^{2}
  =
  \|\psi\|_{\rm pred}^{2}
  +
  \|\phi\|_{\rm pred}^{2}.
\]
Additivity of weights therefore gives
\begin{equation}
\label{eq:cauchy-additive}
  G(x+y)=G(x)+G(y),
  \qquad x,y\ge0.
\end{equation}
Independent renewal composition gives
\begin{equation}
\label{eq:cauchy-multiplicative}
  G(xy)=G(x)G(y).
\end{equation}
Positivity and weak regularity exclude the pathological Cauchy solutions and force \(G(x)=cx\) on \(\R_{\ge0}\).  Normalizing \(W\) to assign weight one to a unit predictive amplitude gives \(c=1\).  Thus \(W(\psi)=\|\psi\|_{\rm pred}^{2}\).  Applying this to the orthogonal sector amplitudes \(\Pi_i\psi\) gives \eqref{eq:born-disintegration}; the mixed-state formula follows by convex extension and gives the usual L\"uders conditional state.
\end{proof}

This result is not a Gleason-type theorem in disguise, because the input is not a noncontextual measure on all projectors.  The input is the renewal structure of stable records: refinement stability, independent renewal composition, and invariance under transformations that are invisible in the predictive quotient.  It is nevertheless compatible with the standard operator-algebraic uniqueness results for quantum probabilities \cite{Gleason1957,Busch2003}.

\subsection{Reset-design closure for implemented bases}
\label{subsec:reset-design}

The remaining question is operational rather than probabilistic.  Once the quadratic weight \eqref{eq:quadratic-weight} is fixed, an arbitrary measurement basis is covered whenever the corresponding instrument can be implemented by the regenerative reset ensemble.  Since probabilities are quadratic in amplitudes, second moments suffice: projective two-design closure is the relevant completeness condition.

Let \(\mathcal E_L\subset U(\Hilb_{\rm pred})\) be the ensemble generated by reset words of length at most \(L\), pushed through the predictive quotient.  Its second frame potential is
\begin{equation}
\label{eq:frame-potential}
  \mathcal F_2(\mathcal E_L)
  =
  \frac{1}{|\mathcal E_L|^2}
  \sum_{U,V\in\mathcal E_L}
  \left|\Tr(U^\ast V)\right|^4 .
\end{equation}
For Haar-random unitaries in dimension at least two the corresponding value is \(2\).  Thus \(\mathcal F_2(\mathcal E_L)-2\) measures the two-design gap \cite{AmbainisEmerson2007,Dankert2009}.

In the finite-cell realization used by the renewal matter and current sectors, the bare real orientation ensemble is not itself a two-design.  The complex predictive phase supplied by the completely positive cell changes the reset closure.  With the characteristic phase
\begin{equation}
\label{eq:cp-phase}
  \delta=\arctan\sqrt5\simeq 65.9^\circ,
\end{equation}
the unphased cell gives a fourth frame moment of order \(7\), whereas the phase-closed reset ensemble approaches the Haar value \(2\) under increasing word length; finite-word computations give values close to \(2.2\) by length \(30\).  The numerical values are not used as axioms in the collapse argument.  Their role is to identify the mechanism by which real orientation data become a sufficiently rich complex reset ensemble for implemented quantum instruments.

\begin{proposition}[Operational Born closure]
\label{prop:operational-born-closure}
Suppose the regenerative reset ensemble is projectively two-design complete in the predictive quotient, or has a two-design gap that vanishes under renewal depth for the class of implemented instruments.  Then the renewal Born disintegration \eqref{eq:born-disintegration} extends from pointer sectors to all measurement bases implementable by the record dynamics.
\end{proposition}

\begin{proof}
Theorem \ref{thm:renewal-born-disintegration} fixes the branch weight as the quadratic predictive norm.  The probabilities of projective and POVM instruments are therefore determined by second moments of the implemented reset ensemble.  A projective two-design reproduces the Haar second moments needed to represent arbitrary implemented basis choices up to the relevant predictive resolution.  Hence any basis generated by the reset-complete instrument family inherits the same quadratic branch weights.
\end{proof}

The status of the Born rule is therefore sharper than a pointer-only statement.  The probability measure is fixed by the renewal disintegration theorem; the finite-cell and reset-design analysis supplies the operational closure condition for measurement bases.  The remaining assumptions are record-structure identifications: the existence of the predictive Hilbert module, stable regenerative central sectors, and reset-design completeness for the instruments under consideration.

\subsection{Pointer basis as renewal stability}

The preferred-basis problem asks why outcomes appear in one basis rather than another.  In renewal language this is not a matter of choosing a basis by convention.  The pointer basis is the stable sector decomposition of the record dynamics.

\begin{definition}[Renewal pointer basis]
A family of projections \(\{P_i\}\) is a renewal pointer basis if repeated application of the record instrument preserves diagonal sectors,
\begin{equation}
\label{eq:pointer-preservation}
  \Instr(P_i\rho P_i)\subseteq P_i\B(\Hilb)P_i
\end{equation}
up to corrections that decay under future regeneration, while off-diagonal sectors are suppressed in the predictive quotient.
\end{definition}

This definition makes the relation to decoherence precise.  Decoherence explains why environmental monitoring suppresses interference in certain bases.  Renewal says that a basis becomes an outcome basis when it is also stable as a regenerative causal-state partition.

\subsection{Objectivity from redundant records}

A record becomes objective when it is not merely available to one observer, but redundantly encoded in many approximately independent fragments of the environment.  This is close to quantum Darwinism and spectrum-broadcast structure \cite{Zurek2009,BrandaoPianiHorodecki2015}.  Renewal interprets redundancy as agreement among subalgebras of one predictive quotient.

\begin{definition}[Redundant renewal record]
A sector label \(i\) is redundantly objective if there exist subalgebras \(\Ralg_1,\ldots,\Ralg_N\) such that each \(\Ralg_k\) contains enough information to infer \(i\), and the conditional future predictions agree on overlaps of the corresponding observer algebras.
\end{definition}

\begin{proposition}[Objectivity criterion]
\label{prop:objectivity}
If a sector label is redundantly encoded in many approximately independent record subalgebras, then observers with access to different fragments condition on the same sector with probability approaching one as redundancy grows.  The outcome is objective within the shared predictive quotient.
\end{proposition}

Objectivity is therefore neither a primitive classical fact nor a private mental event.  It is the fixed point of redundant record agreement.

\section{Relation to the main interpretations}
\label{sec:interpretations}

The renewal account is best classified by what it refuses and what it adds.  It refuses a new objective-collapse force.  It adds a record algebra to unitary/decoherence-based accounts.  The result is many-worlds-compatible, but not identical to the strongest form of Everettian metaphysics.

\subsection{Comparison table}

\begin{center}
\small
\renewcommand{\arraystretch}{1.15}
\begin{tabularx}{\textwidth}{@{}p{0.23\textwidth}X@{}}
\toprule
\textbf{Interpretation} & \textbf{Relation to renewal collapse}\\
\midrule
Copenhagen / operational & Agrees that collapse is an update of predictions, but replaces the primitive classical cut by an emergent regenerative record algebra.\\
Decoherence & Very close dynamically: environmental records suppress interference.  Renewal adds the criterion that effective branches are stable predictive sectors.\\
Many-Worlds & Compatible if global evolution remains unitary.  Renewal weakens the ontology of ``worlds'' and sharpens branch definition through records.\\
Consistent histories & Close in spirit: probabilities attach to decoherent families.  Renewal selects physically relevant histories by regenerative predictive stability.\\
Relational quantum mechanics & Close: state assignments are relative.  Renewal specifies the algebraic record structure relative to which they are assigned.\\
QBism & Shares the rejection of collapse as a physical shockwave, but treats records as objective algebraic structures rather than purely agent-personal commitments.\\
GRW/CSL objective collapse & Largely opposed: no stochastic localization term is added, and no generic GRW/CSL deviations are predicted.\\
Bohmian mechanics & Different ontology: the primitive variables are record sectors and predictive channels, not hidden particle configurations guided by a pilot wave.\\
\bottomrule
\end{tabularx}
\end{center}

The table should not be read as an attempt to collapse all interpretations into one.  Rather, it locates renewal collapse in the existing landscape.  Its natural neighbors are decoherence, consistent histories, relational interpretations, and Everett-compatible unitarity.  Its distinctive ingredient is the predictive quotient.

\subsection{Weak branch realism}

The question most likely to arise is whether the renewal view is simply Many-Worlds.  It is not, although it can be embedded in an Everettian ontology.  The Everettian claim is that the universal state never collapses and that all decohered branches are physically real.  Renewal theory can accept global unitarity, but it does not need to assign equal primitive status to every decohered component of every decomposition.  A branch becomes physically meaningful when it is a stable sector of a record algebra.

\begin{interpretation}[Weak branch realism]
\label{interp:weak-branch}
Branches are real as stable predictive sectors of regenerative record algebras.  This is weaker than strong world realism, according to which every decohered component of the universal wavefunction is an equally fundamental world, and stronger than a purely instrumentalist view in which branches are only bookkeeping devices.
\end{interpretation}

Weak branch realism has a technical advantage.  It avoids treating arbitrary basis decompositions as worlds.  The branch decomposition is fixed only after the record algebra and its regenerative stability have been specified.

\subsection{Why this is not objective collapse}

Objective-collapse models modify quantum dynamics so that macroscopic superpositions are physically unstable.  Renewal theory makes a different prediction.  If one searches for universal localization noise of the GRW/CSL type, the renewal account does not require it.  Apparent collapse is produced by conditioning on records, not by a new random field.  Therefore the most direct experimental deviations expected in objective-collapse theories are not generic predictions of the renewal interpretation.

This does not make renewal collapse empirically empty.  It shifts the possible empirical signatures toward record dynamics: decoherence rates, redundancy scaling, reset-design statistics, measurement-induced transitions, non-Markovian record correlations, and, in the horizon setting, deviations from exactly memoryless radiation.  The object of study is the record stream, not an independent collapse force.

\section{Phenomenology of record streams}
\label{sec:phenomenology}

The renewal account is experimentally conservative with respect to objective-collapse searches and more specific about measurement records.  It does not predict a universal GRW/CSL localization noise, a mass-proportional stochastic field, or a fundamental violation of Schr\"odinger evolution.  Its empirical content lies instead in the statistics of record formation and record propagation.

The relevant observables are those that test whether a stable sector has formed in the predictive quotient.  Four classes are especially natural.

\begin{center}
\small
\renewcommand{\arraystretch}{1.14}
\begin{tabularx}{\textwidth}{@{}p{0.26\textwidth}X@{}}
\toprule
\textbf{Experimental regime} & \textbf{Renewal target}\\
\midrule
Decoherence and recoherence tests & Determine whether off-sector coherences are merely small in a chosen environment or actually invisible to the retained predictive algebra.\\
Quantum Darwinism and spectrum-broadcast tests & Measure redundancy growth and agreement among environmental fragments as a proxy for objectivity of the record sector.\\
Reset and random-circuit experiments & Estimate second frame moments and two-design gaps of the implemented reset ensemble; the Born-disintegration argument predicts that operational basis closure is controlled by these second moments.\\
Measurement-induced transitions & Track changes in the regenerative record algebra as monitoring strength, feedback, or reset depth is varied.\\
\bottomrule
\end{tabularx}
\end{center}

In this sense the renewal phenomenology is closer to decoherence, quantum Darwinism, and random-circuit measurement physics than to dynamical-reduction phenomenology.  The decisive question is not whether isolated systems undergo spontaneous localization in the absence of records.  It is whether the record stream becomes stable, redundant, and reset-complete enough to support the conditional probabilities assigned by \cref{thm:renewal-born-disintegration,prop:operational-born-closure}.

The same distinction clarifies the role of black holes.  A perfectly thermal, exactly memoryless radiation stream would fail as a complete record, whereas a non-Markovian radiation stream with late-time purification can in principle become record-complete.  The horizon discussion below is therefore not an additional collapse mechanism.  It is the most severe version of the same record-stream test.

\section{The horizon record test}
\label{sec:horizon}

Black holes enter the present paper only through the information-theoretic role of records.  The detailed horizon dynamics is assumed from the renewal black-hole construction.  What matters here is the following translation: Hawking radiation is an exterior record stream, and the Page curve is a statement about the completeness of that stream.

In ordinary measurement, the record is produced by an apparatus and its environment.  The problem is branch selection: why does an observer experience one outcome?  In black-hole evaporation, the exterior observer receives radiation from a horizon.  The problem is branch reconstruction: does the radiation contain enough correlations to recover the information about what entered the black hole?  These are dual uses of the same notion of record.

\begin{center}
\small
\renewcommand{\arraystretch}{1.15}
\begin{tabularx}{\textwidth}{@{}p{0.30\textwidth}p{0.31\textwidth}X@{}}
\toprule
& \textbf{Laboratory measurement} & \textbf{Black-hole evaporation}\\
\midrule
Record carrier & apparatus/environment fragments & Hawking radiation and horizon-accessible exterior algebra\\
Central question & branch selection & branch reconstruction\\
Effective collapse & conditioning on a pointer sector & conditioning on exterior radiation records\\
Failure mode & no stable outcome basis & exactly thermal memoryless radiation\\
Success condition & redundant record objectivity & Page-complete radiation record\\
\bottomrule
\end{tabularx}
\end{center}

\begin{principle}[Horizon record principle]
\label{principle:horizon-record}
The black-hole information problem is the horizon analogue of the measurement record problem.  The exterior algebra must receive a record stream whose correlations become complete enough to purify the radiation after the Page time.  If the stream were exactly thermal and memoryless, it would fail as a complete regenerative record.
\end{principle}

This principle does not by itself prove unitarity, derive an island formula, or compute a radiation density matrix.  It gives the interpretive bridge: the same algebraic condition that makes a laboratory outcome objective is, at horizon scale, the condition that makes exterior reconstruction possible.

\begin{proposition}[Selection-reconstruction duality]
\label{prop:selection-reconstruction}
Let \(\Ralg\) be a regenerative record algebra.  In a measurement setting, conditioning on a stable central sector of \(\Ralg\) produces branch selection.  In a horizon setting, increasing access to a sufficiently complete radiation subalgebra produces branch reconstruction.  Both processes are governed by inclusion and conditioning of record algebras.
\end{proposition}

\begin{proof}
In measurement, the observer begins with access to the apparatus/environment sector label and conditions the future state on that label.  In the horizon problem, the exterior observer initially has an incomplete radiation algebra; as additional radiation is included, correlations with the remaining black-hole degrees of freedom determine whether the exterior algebra becomes complete.  In both cases, the operative maps are restriction to an accessible algebra, enlargement of the record algebra, and conditionalization on its stable sectors.
\end{proof}

This is the proper scope of the black-hole discussion in a collapse paper.  The horizon is not invoked to generate collapse.  It is invoked because it tests the same claim under maximal stress: whether records, and records alone, can carry the information that makes a branch definite or recoverable.

\section{Conclusion and outlook}
\label{sec:outlook}

The renewal interpretation of wavefunction collapse can be summarized in one sentence: collapse is the effective projection onto a stable predictive record sector.  The proposal is therefore neither a new collapse dynamics nor a purely subjective update.  It is a record-relative account of definiteness.  The state that matters for an observer is the state restricted to the algebra of records that survive future predictive equivalence.

The resulting position has a definite place among quantum interpretations.  It is close to decoherence because interference is suppressed by environmental records.  It is close to consistent histories because probabilities are meaningful only for stable histories or sectors.  It is close to relational views because state assignments are relative to an accessible record algebra.  It is compatible with Everettian unitarity because the global state need not collapse.  But it is not simply Many-Worlds: branch structure is not arbitrary, and the ontology of branches is weakened to stable predictive sectors.  It is also not GRW/CSL: no universal stochastic localization field is added.

The strongest closed claims are the following.
\begin{center}
\small
\renewcommand{\arraystretch}{1.12}
\begin{tabularx}{\textwidth}{@{}p{0.34\textwidth}p{0.20\textwidth}X@{}}
\toprule
\textbf{Claim} & \textbf{Status} & \textbf{Comment}\\
\midrule
No fundamental collapse force & Structural stance & Renewal collapse is not objective-collapse dynamics.\\
Collapse as record conditioning & Core thesis & The measurement update is conditionalization on a stable record sector.\\
Pointer basis from renewal stability & Strong conditional & Pointer states are stable causal states of the record instrument.\\
Pointer-sector Born weights & Theorem & Follows from additivity on the Boolean algebra of stable records.\\
Quadratic branch weights & Theorem & Forced by predictive refinement additivity, independent renewal multiplicativity, and phase/deck invariance.\\
Full operational Born rule & Conditional theorem & Holds for all measurement bases implementable by a reset ensemble with projective two-design closure.\\
Objectivity & Strong conditional & Redundant records produce intersubjective agreement.\\
No-signalling & Standard theorem & Collapse is algebra-relative conditioning, not a physical signal.\\
Record-stream phenomenology & Concrete target & Tests concern redundancy growth, reset-design gaps, non-Markovian correlations, and measurement-induced transitions.\\
Black-hole information link & Interpretive application & Page recovery is record completeness in the horizon setting.\\
\bottomrule
\end{tabularx}
\end{center}

Several problems remain.  The first is to compute renewal einselection rates from explicit record instruments rather than treating stability abstractly.  The second is to prove reset-design completeness analytically for the full class of experimentally implementable instruments, beyond finite-word and finite-sector computations.  The third is to characterize measurement-induced transitions in the interacting renewal regime.  The fourth is to make the horizon record-completeness criterion precise enough to compare with black-hole Page dynamics, without importing more black-hole machinery than the measurement problem requires.

The conceptual conclusion is nevertheless already sharp.  The renewal programme does not attempt to explain measurement by adding a second dynamics.  It explains measurement by identifying the algebraic event that makes a quantum alternative predictive: the formation of a stable regenerative record.  Collapse is the observer's passage into such a sector; objectivity is the redundancy of that sector; and black-hole recovery is the demand that an exterior radiation stream eventually supply a complete enough sector algebra to reconstruct what the horizon concealed.

\begin{thebibliography}{99}

\bibitem{PelissierLorentzian}
A.~P\'elissier,
\emph{Renewal Spectral Geometry and the Emergence of Lorentzian Spacetime},
arXiv placeholder.

\bibitem{PelissierEinstein}
A.~P\'elissier,
\emph{Emergent Einstein Dynamics, de Sitter Horizons, and the Quantum-Gravity Layer from Renewal Memory},
arXiv placeholder.

\bibitem{PelissierBridge}
A.~P\'elissier,
\emph{Renewal ER=EPR: Predictive Screens, Deck-Odd Correlation, and Two-Sided Bridges},
arXiv placeholder.

\bibitem{PelissierK4}
A.~P\'elissier,
\emph{Renewal Currents and Cycle Geometry},
arXiv placeholder.

\bibitem{PelissierBorn}
A.~P\'elissier,
\emph{Renewal Born Disintegration and Reset Designs},
arXiv placeholder.

\bibitem{PelissierBlackHoles}
A.~P\'elissier,
\emph{Renewal Black Holes: Horizon Response, Non-Markovian Evaporation, Chiral Horizon Noise, and the Trans-Renewal Interior},
arXiv placeholder.

\bibitem{Gleason1957}
A.~M.~Gleason,
\emph{Measures on the closed subspaces of a Hilbert space},
J. Math. Mech. \textbf{6} (1957), 885--893.

\bibitem{Busch2003}
P.~Busch,
\emph{Quantum states and generalized observables: a simple proof of Gleason's theorem},
Phys. Rev. Lett. \textbf{91} (2003), 120403.

\bibitem{AmbainisEmerson2007}
A.~Ambainis and J.~Emerson,
\emph{Quantum \(t\)-designs: \(t\)-wise independence in the quantum world},
Proc. 22nd IEEE Conference on Computational Complexity (2007), 129--140.

\bibitem{Dankert2009}
C.~Dankert, R.~Cleve, J.~Emerson, and E.~Livine,
\emph{Exact and approximate unitary 2-designs and their application to fidelity estimation},
Phys. Rev. A \textbf{80} (2009), 012304.

\bibitem{Everett1957}
H.~Everett,
\emph{Relative state formulation of quantum mechanics},
Rev. Mod. Phys. \textbf{29} (1957), 454--462.

\bibitem{DeWittGraham1973}
B.~S.~DeWitt and N.~Graham (eds.),
\emph{The Many-Worlds Interpretation of Quantum Mechanics},
Princeton University Press, 1973.

\bibitem{Wallace2012}
D.~Wallace,
\emph{The Emergent Multiverse: Quantum Theory according to the Everett Interpretation},
Oxford University Press, 2012.

\bibitem{Zeh1970}
H.~D.~Zeh,
\emph{On the interpretation of measurement in quantum theory},
Found. Phys. \textbf{1} (1970), 69--76.

\bibitem{Zurek1981}
W.~H.~Zurek,
\emph{Pointer basis of quantum apparatus: into what mixture does the wave packet collapse?},
Phys. Rev. D \textbf{24} (1981), 1516--1525.

\bibitem{JoosZeh1985}
E.~Joos and H.~D.~Zeh,
\emph{The emergence of classical properties through interaction with the environment},
Z. Phys. B \textbf{59} (1985), 223--243.

\bibitem{Schlosshauer2007}
M.~Schlosshauer,
\emph{Decoherence and the Quantum-to-Classical Transition},
Springer, 2007.

\bibitem{Zurek2009}
W.~H.~Zurek,
\emph{Quantum Darwinism},
Nature Phys. \textbf{5} (2009), 181--188.

\bibitem{BrandaoPianiHorodecki2015}
F.~G.~S.~L.~Brand\~ao, M.~Piani, and P.~Horodecki,
\emph{Generic emergence of classical features in quantum Darwinism},
Nature Commun. \textbf{6} (2015), 7908.

\bibitem{Griffiths1984}
R.~B.~Griffiths,
\emph{Consistent histories and the interpretation of quantum mechanics},
J. Stat. Phys. \textbf{36} (1984), 219--272.

\bibitem{Omnes1992}
R.~Omn\`es,
\emph{Consistent interpretations of quantum mechanics},
Rev. Mod. Phys. \textbf{64} (1992), 339--382.

\bibitem{GellMannHartle1993}
M.~Gell-Mann and J.~B.~Hartle,
\emph{Classical equations for quantum systems},
Phys. Rev. D \textbf{47} (1993), 3345--3382.

\bibitem{Rovelli1996}
C.~Rovelli,
\emph{Relational quantum mechanics},
Int. J. Theor. Phys. \textbf{35} (1996), 1637--1678.

\bibitem{FuchsMerminSchack2014}
C.~A.~Fuchs, N.~D.~Mermin, and R.~Schack,
\emph{An introduction to QBism with an application to the locality of quantum mechanics},
Am. J. Phys. \textbf{82} (2014), 749--754.

\bibitem{GRW1986}
G.~C.~Ghirardi, A.~Rimini, and T.~Weber,
\emph{Unified dynamics for microscopic and macroscopic systems},
Phys. Rev. D \textbf{34} (1986), 470--491.

\bibitem{Pearle1989}
P.~Pearle,
\emph{Combining stochastic dynamical state-vector reduction with spontaneous localization},
Phys. Rev. A \textbf{39} (1989), 2277--2289.

\bibitem{BassiGhirardi2003}
A.~Bassi and G.~C.~Ghirardi,
\emph{Dynamical reduction models},
Phys. Rept. \textbf{379} (2003), 257--426.

\bibitem{Bell1964}
J.~S.~Bell,
\emph{On the Einstein Podolsky Rosen paradox},
Physics \textbf{1} (1964), 195--200.

\bibitem{Page1993}
D.~N.~Page,
\emph{Information in black hole radiation},
Phys. Rev. Lett. \textbf{71} (1993), 3743--3746.

\bibitem{HaydenPreskill2007}
P.~Hayden and J.~Preskill,
\emph{Black holes as mirrors: quantum information in random subsystems},
JHEP \textbf{09} (2007), 120.

\end{thebibliography}

\end{document}
